\documentclass[10pt,a4paper]{article}
\usepackage{amsmath,amssymb,amsthm,graphicx}
\usepackage{hyperref}
\usepackage{latexsym}
\usepackage{amssymb}
\usepackage{amsthm}
\usepackage{ulem}
\usepackage{enumitem}
\setlist{itemsep=0.4pt,topsep=0pt}
\parindent 0pt
\parindent 0pt
\parskip 12pt
\newcommand{\HCF}{\mathrm{hcf}}
\newcommand{\ZZ}{\mathbb{Z}}
\newcommand{\rk}{\mathrm{rk}}
\DeclareMathOperator{\len}{length}
\DeclareMathOperator{\bl}{b}
\DeclareMathOperator{\cl}{cl}
\DeclareMathOperator{\di}{d}
\DeclareMathOperator{\Cl}{Cl}
\DeclareMathOperator{\acl}{acl}
\DeclareMathOperator{\tp}{tp}
\DeclareMathOperator{\Th}{Th}
\newcommand{\Cor}{\mathcal{C}_\mathrm{or}}
\newcommand{\Cun}{\mathcal{C}_\mathrm{un}}
\newcommand{\mons}{\mathbb{M}}
\newcommand{\strong}{\mathrm{strong}}
\newcommand{\axioms}{\Delta_{\mathrm{FSTS}}}
\newtheoremstyle{note}% hnamei
{8pt}% hSpace abovei
{14pt}% hSpace belowi
{\normalfont}% hBody fonti
{}% hIndent amounti1
{\bfseries}% hTheorem head fonti
{ }% hPunctuation after theorem headi
{.5em}% hSpace after theorem headi2
{}% hTheorem head spec (can be left empty, meaning ‘normal’)i
\theoremstyle{note}
\newtheorem{thm}{Theorem}[section]
%\newtheorem{theorem}{Theorem}
\newtheorem{lemma}[thm]{Lemma} 
\newtheorem{df}[thm]{Definition}
\newtheorem{rmk}[thm]{Remark} 
\newtheorem{prop}[thm]{Proposition}
\newtheorem{fact}[thm]{Fact}
\newtheorem{cor}[thm]{Corollary}
\newtheorem{conj}[thm]{Conjecture}
\newtheorem{ex}[thm]{Example}
\newtheorem{qu}[thm]{Question}
\newtheorem*{notation}{Notation} 
\title{Above every element there is a binary tree of arbitrary finite height}
\author{Silvia Barbina, Enrique Casanovas}
\date{arxiv:2411.13723v3; Creative Commons Attribution 4.0 International}
\begin{document}
\maketitle
\noindent\textit{Editorial note.} The counted statement is the author's environment at source0.tex lines 841-846 and its proof environment opens at line 847; the pairing is by adjacency and the retained proof range is the whole environment.\par\medskip
\noindent\textit{Editorial note.} The preamble is the author's own package lines, macro block and theorem-environment declarations (source lines 4-65); the author's optioned amsmath line and the fleqn-dependent layout line are not retained.\par\medskip
\noindent\textit{Editorial note.} The source carries no label-number dump and no counter values are imposed, so the printed numbering is the wrapper's default; identity is the author label, or the pinned source position where the author gave none.\par\medskip
\noindent\textit{Editorial note.} Metadata provenance: the source record of the pinned copy and the licence evidence file are both present; the licence is CC BY 4.0 with status supported, and the card's source\_version matches the value recorded in that evidence file.\par\medskip
\noindent\textit{Editorial scope.} One original result of the pinned author manuscript, counted once: binary trees of arbitrary finite height above an element (author label tree3).\par\medskip
\noindent\textit{Allowed standard prerequisites (editorial).}\begin{itemize}
\item Binary tree in a Steiner triple system. A family indexed by finite binary sequences whose nodes are ordered by extension and whose elements grow with the index.. It is the object constructed.
\item Products above a prescribed element. Below every decomposition of an element there are two elements above any prescribed bound.. It is the fact that lets the construction continue at each node.
\end{itemize}\medskip
\begin{lemma}\label{tree2}
	\begin{enumerate}
		\item An HF-ordering of an infinite STS does not have a maximum.
		\item If $<$ is an HF-ordering of an infinite Steiner triple system $M$, then for every $a\in M$, for every  $b \geq a$, there are $b_1,b_2>b$ such that  $a=b_1\cdot b_2$.
	\end{enumerate}
\end{lemma}

\begin{prop}\label{tree3} Let $<$ be an HF-ordering of the infinite Steiner triple system $M$ and let $n<\omega$. For every $a\in M$    there is a binary tree $T=\{a_s\mid s\in 2^{<n}\}$  in $M$  such that
	\begin{enumerate}
		\item $a_s>a$  for every $s$
		\item $a_s>a_t$  if  $s\subsetneq t$.
	\end{enumerate}
\end{prop}

\begin{proof}
%	By induction on $n$. The initial cases $n=0,1$ are trivial. Consider the case  $n=2$. Let $b>a$. We use Lemma~\ref{tree2} to get some $b_1 >b_2 >b$ in $M$ such that $b=b_1\cdot b_2$. We put $a_\emptyset = b_1$,  $a_0=b_2$ and $a_1 = b$. 
 

%	Now we assume inductively the result holds for $n\geq 2$.
%	Given $a\in M$, the inductive hypothesis gives some tree $T_0=\{a_s\mid s\in 2^{<n}\}$ as required for $a$.  Again by the inductive hypothesis we can find in $M$ a tree $T_1=\{a^\prime_s\mid s\in 2^{<n}\}$ with $a^\prime_s>a_\emptyset$ and condition 2 holds. Now we combine $T_0$ and $T_1$ to get a binary tree $T=\{b_s \mid s\in 2^{<n+1}\}$. We set $b_{0s}= a_s$, $b_{1s}= a^\prime_s$  and  $b_\emptyset = a_\emptyset\cdot a^\prime_\emptyset$.  It is enough to check that   $a_\emptyset\cdot a^\prime_\emptyset > a^\prime_\emptyset$. Since $n\geq 2$ this is clear, since otherwise the blocks $\{a^\prime_\emptyset,a^\prime_0,a^\prime_1\}$  and  $\{a^\prime_\emptyset, a_\emptyset, a_\emptyset \cdot a^\prime_\emptyset\}$ are different and contradict that $<$ is an HF-ordering.
	We prove by induction on $n$ that  for every $a\in M$ for every finite $A\subseteq M$, there is a binary tree $T=\{a_s\mid s\in 2^{<n}\}$  in $M$  satisfying conditions 1,2 and such that there is no block connecting elements of $T$ with elements of $A$ (that is, no products $t \cdot a$ with $t \in T$ and $a \in A$ are in $T \cup A$).
	 The initial cases $n=0,1$ are trivial. Consider the case  $n=2$. Let $A^\prime$ be the set of products $\{x\cdot y\mid x,y\in A\}$, and let $b>a$ and larger than every element of $A^\prime$. We use Lemma~\ref{tree2} to get some $b_1 >b_2 >b$ in $M$ such that $b=b_1\cdot b_2$. We let $a_\emptyset = b_1$,  $a_0=b_2$ and $a_1 = b$. 
	
	
	Now we assume inductively the result holds for $n\geq 2$.
	Given $a\in M$ and a finite set $A\subseteq M$, the inductive hypothesis gives some tree $T_0=\{a_s\mid s\in 2^{<n}\}$ as required for $a,A$.  Let $A^\prime =A\cup T_0$. Again by inductive hypothesis, we can find a tree $T_1=\{a^\prime_s\mid s\in 2^{<n}\}$ in $M$  with $a^\prime_s$ larger than every element of $A^\prime$, with no blocks connecting $T_1$ and $A^\prime$  and such that condition 2 holds. Now we combine $T_0$ and $T_1$ to get a binary tree $T=\{b_s \mid s\in 2^{<n+1}\}$. We set $b_{0s}= a_s$, $b_{1s}= a^\prime_s$  and  $b_\emptyset = a_\emptyset\cdot a^\prime_\emptyset$.  We first  check that   $a_\emptyset\cdot a^\prime_\emptyset > a^\prime_\emptyset$. Since $n\geq 2$ this is clear, since otherwise the blocks $\{a^\prime_\emptyset,a^\prime_0,a^\prime_1\}$  and  $\{a^\prime_\emptyset, a_\emptyset, a_\emptyset \cdot a^\prime_\emptyset\}$ are different and contradict that $<$ is an HF-ordering. Since there is no block connecting elements of $T_1$ with elements of $T_0$, $T$ is a tree.
\end{proof}

\end{document}
